\section{The Living Sacrifice}

Paul's ``therefore'' joins the denial that anyone first gives to God, the confession that all things are from him and through him and for him, and the mercy that precedes response: the creature is to offer the body as a living sacrifice (Rom 11:35--12:1). The offering neither repairs a divine lack nor transfers an independently possessed good into God's estate. The body offered is already created; the reason renewed is already sustained; the will that consents already acts through powers it has received.

Sacrifice makes the return visible. Aquinas argues that natural reason recognizes some outward sign of the honor and subjection due to the first principle, while positive law determines its concrete forms. Embodied signs subject the mind to God and can arouse the devotion they express; the outward offering signifies the inward sacrifice in which the soul offers itself to God (ST II--II, q.~81, a.~7; q.~85, aa.~1--2). The sign is false when it is severed from the inward act, but the inward act is not therefore disembodied. An embodied and social creature answers with more than an intention hidden from the world.

For Augustine, every work truly ordered to holy fellowship with God is sacrifice; the person becomes an offering, works of mercy are referred to the supreme good, and the whole redeemed city is offered through the High Priest who first offered himself (\emph{City of God} X.5--6, 20). Christian self-offering is not an autonomous ascent. It occurs within Christ's prior oblation.

The Eucharist is not deduced from First-Mover metaphysics. Christ's sacrifice is a genuinely human oblation offered by the divine Son; it cannot be placed alongside creaturely returns as though commensurable with them. Christ is priest and victim; the Church's self-offering depends wholly upon his perfect prior oblation. United to him, she offers the divine Victim and learns to offer herself with him (\emph{Sacrosanctum Concilium} 47--48; \emph{Lumen Gentium} 10--11). Work and prayer, family life and rest, hardship and acts of mercy become spiritual sacrifices when offered through Christ and brought into his Eucharistic offering (\emph{Lumen Gentium} 34).

A living sacrifice takes the received self as its matter: intellect, will, body, time, possessions, and common life. It does not merely acknowledge that creaturely being is from God. It places creaturely action within that truth.

Because Christian sacrifice is ecclesial, its words, signs, times, and offices are received rather than privately devised. Authority belongs inside the account of worship; its presence does not by itself convert command into coercion.
